% a06_BALtorus_d7b: continuation of a06_BALtorus_d7 (input from sec/a06_BALtorus_c.tex): the regrouped local coefficients, the induction of pp. 444-445, the assembled Proposition 1 at L = 5, and the one open row.

\begin{theorem}[bounds on the regrouped local coefficients]\label{thm:a06d:S48}
For every useful dangerous node $S$ of type $2A$ or $\phi\phi A$ the regrouped local coefficients satisfy: $2A$: $|c^{(0)}(\hat x)|\le Cs(v_S)^{D(S)}$ (zeroth moment), and for $D(S)=-1$ also $|c^{(1)}_\mu|\le C$ (first moment); $\phi\phi A$: $|c^{(0)}|\le C$; with $C$ independent of $\eta,k,\eps$, the labels, $\Omega$, $\tilde B$ and the torus. For $2\phi$ nodes the zeroth order cancels exactly against the counterterm member, and the $D=-1$ first moment has the bounded regrouped coefficient of \cite[(3.11)]{Ba3}.
\end{theorem}

\begin{proof}
Induction on the number of lines of the node. \emph{The value.} By \cref{lem:a06d:S41}(a), $c$ is the regrouped value; rescaled to $\xi=L^{-j_0}$ it is $s(v_S)^{D(S)}$ times the same expression on $\xi\Z^3$, with lines $\breve G_l=C^\xi+R_l$; smooth multipliers are moved to $\hat x$ (\cref{lem:a06d:S27}, gain $1$), sharp ones are treated by \cref{lem:a06d:S46}. \emph{The all-$C^\xi$ part} is the per-label full class with full lines $C^\xi$, internal sums unrestricted on $\xi\Z^3$, and every counterterm member re-expanded with its own lines (\cite[p.~445]{Ba3}: ``the same transformations in mass renormalization counterterms''), inner counterterm coefficients being fixed numbers, which \cref{lem:a06d:S26} allows: the $2A$ zeroth moment is $0$ by \cref{lem:a06d:S26}; the $2A$ first moment and the $\phi\phi A$ zeroth moment are $0$ by the reflection identities of \cite[(3.30)]{Ba3}, applicable because the coefficients are now constant; the torus-versus-$\xi\Z^3$ identification is the image-term argument of \cite[pp.~418--419]{Ba3}. \emph{Terms with at least one $R$-line, $D(S)=-1$}: these are one-loop nodes, bounded by the cancellation $\sum_{y'}(\partial C\partial^*)(y'-y)=0$ for the term $R\cdot\partial C\partial^*$, which uses only the bounds on $R$ and the forward-difference structure of the other line, so it applies also to $R'$-lines (\cref{lem:a06d:S47}(c)) and to slices on which the derivative acts. \emph{Terms with at least one $R$-line, $D(S)=0$}: re-slice every line (\cref{lem:a06d:S47}(a), (b)) and apply the bound function estimate \cite[(2.16)]{Ba3} to the coefficient (anchor fixed, external fields factored out). Along each ordering: a component containing an $R$-slice has degree $\ge D(H)+1$, which is $\ge\frac12$ unless $D(H)=-1$, and a $D(H)=-1$ component is a one-loop $2\phi$ or $2A$ node with an $R$-line, processed by the previous case ($2A$) or by its counterterm member with the same $R$-line ($2\phi$, exact; the first moment bounded as for \cite[(3.11)]{Ba3}); a component without an $R$-slice has its original degree, and the dangerous ones are inner nodes, processed by the induction hypothesis (local coefficients bounded, remainders gaining $\alpha$) and by \cref{lem:a06d:S44}; the whole coefficient has degree $\ge D(S)+1=1$; inner $2\phi$ nodes whose counterterm assignments are useless relative to the inner scale are bounded by \cref{lem:a06d:S45} (loss $\alpha'$). So every component has degree $\ge\min\{\frac12,\alpha\}-\alpha'>0$, and \cite[(2.16)]{Ba3} bounds the sum uniformly in $\xi$. \emph{$\phi\phi A$ remainder terms}: one-loop, so the inner components are trees, the only dangerous ones \cite[(2.4)]{Ba3}, handled by integration by parts inside the coefficient; the previous case applies with no inner nodes.
\end{proof}

\begin{theorem}[{\cite[(3.5)]{Ba3} for every class}]\label{thm:a06d:S49}
For every class component after the cube steps, $\sum_G\sum_{\tilde l}\sum_{j\in J(\tilde l)}E(G(j))$ is a finite sum of terms each of which is either a generalised bound function in which every component along every ordering of its support has degree $\ge\theta\defeq\alpha/2$, or is cancelled exactly (vacuum components, $2\phi$ zeroth orders). Hence \cite[(3.3)]{Ba3} holds by the bound function estimate \cite[(2.16)]{Ba3}.
\end{theorem}

\begin{proof}
Fix an ordering, i.e.\ an index assignment class $J(\tilde l)$, and process the dangerous nodes along $\tilde l$ innermost first. \emph{Order and anchors}: non-\cite[(2.4)]{Ba3} nodes innermost first with any anchors ($2\phi$: the counterterm anchor); each outer node's split applied to all terms produced by the inner processing, $\Delta$-terms included (\cref{lem:a06d:S42}); \cite[(2.4)]{Ba3} components treated inside the smallest dangerous non-\cite[(2.4)]{Ba3} node containing them (\cref{lem:a06d:S44}), or alone by \cite[(2.8)]{Ba3}. \emph{Local terms}: $2\phi$, the zeroth order cancels exactly against the counterterm member at the same indices in the useful assignments, and the $D=-1$ first moment has its bounded coefficient; $2A$ and $\phi\phi A$, the local vertex $v_S$ has coefficient bounded by \cref{thm:a06d:S48}, its value being the regrouped one (\cref{lem:a06d:S41}(a)); for $2A$ with $D=-1$ this vertex is the non-vanishing zero-momentum remainder vertex of \cite[(3.30)]{Ba3}, of power $-1$ at $s(v_S)$, i.e.\ degree $2+D(S)=1$ in the convention \cite[(2.3)]{Ba3}. \emph{Remainder terms} gain $\alpha$ or $1+\alpha$ inside $S$ and lose the same on $S$'s external lines (\cref{lem:a06d:S41}(b)); where an inner processing acts on an outer node's external field, the outer node gains at least $\alpha$ and has no local part in that term (\cref{lem:a06d:S42}); so after all nodes are processed every component that was dangerous has degree $\ge\alpha$, and every other component keeps its degree, which is $\ge\frac12$ since degrees lie in $\frac12\Z$. \emph{Useless counterterm assignments} are bounded by \cref{lem:a06d:S45} at a loss $\alpha'=\alpha/2$. \emph{Remainder vertices inside larger graphs}: a remainder vertex of a $2A$ node $S$ has power $D(S)$ at $s(v_S)$; components containing it have the degree of the same component with $S$ restored (\cref{lem:a06d:S41}(b)); if such a component is itself dangerous it is a node of the next level and is processed there, its local coefficient being the regrouped value with $S$ unprocessed (\cref{lem:a06d:S41}(a)), so the remainder vertex never needs a bound beyond its power. \emph{Induction}: this is Ba{\l}aban's ``for each ordering \dots\ the last divergent subgraph \dots\ apply (3.5) \dots\ then apply (3.19), (3.31) or (3.34) to $G_0$ itself'' \cite[pp.~444--445]{Ba3}, an induction on the number of lines, \cref{thm:a06d:S48} being applied inside it to coefficients of strictly smaller nodes. \emph{Vacuum and exception classes}: vacuum components cancel against \cite[(1.24)]{Ba3} up to image terms, and $1A$/$3A$ factors vanish at $B=0$.
\end{proof}

\paragraph{\cite[Prop.~1, (1.33)]{Ba3} at $L=5$.} With the steps above, every one of the twenty sketched steps and seven further places of \cite[pp.~420--445]{Ba3} is a written proof: the external-field readings of \cite[(1.16)]{Ba3} and of $\delta G_k(\Omega,\Omega_2,\tilde B)$ \cite[pp.~414, 424]{Ba3} (the latter at depth, as (I2a) of Step~3(D$_2$)), the slices and the summation by parts \cite[(2.6)--(2.12)]{Ba3}, the induction \cite[(2.16)]{Ba3}, \cite[Prop.~2.2, (2.17), Cor.~2.3]{Ba3}, the Ward--Takahashi identities for full lines and per label (\cref{lem:a06d:S26}), \cite[(3.1)]{Ba3} in the cluster form (\cref{cor:a06d:S24}), the regrouping (\cref{thm:a06d:S23}), the component cube of side $l_N$, the cube split with $l^\sharp$ (\cref{thm:a06:BALA}(a)), the constant field $B_0$ with the $k$-uniform constant of Step~3(C), the $\tilde B'$-expansion and its classes, the chains without old vertex (by an $L^2$ chain bound), gauge covariance and the gauge-out of $B_0$, $\delta G_k(\square,\eta\Z^d,0)$, the norms \cite[(1.32)]{Ba3} at $B_0$ and at $\tilde B$ (the constant carrying $(1+\operatorname{diam}\supp f)^{1-\alpha}$), and the renormalisation (\cref{thm:a06d:S49}). The inputs outside the proposition are the propagator estimates of group (A) (\cref{thm:a06A:groupA}) and \cref{thm:a06:A5}; the parameter rows, all fixed after $L$, $\bar n$, $R_0$ and before $\eps_0$, are $R\ge R_0+5$ (the depth of \cite[(3.1)]{Ba3}), $R\ge R_{\mathrm{flat}}+R_{\mathrm{Ba}}+2$ ($R_{\mathrm{flat}}$ the depth parameter of the cube step, $R_{\mathrm{Ba}}\ge R_0$ the depth of the H\"older clause of the background, $R_0$ the constant defined in \cite[(2.19)]{Ba4}; the H\"older depth of the cube step), $l^\sharp=l_N+4V^\sharp_{\mathrm{lab}}(\bar n)$ in \cref{thm:a06:BALA}(a), torus side $\ge\max\{6\sqrt3\,r(t_k)+2,8/\delta_Z\}$, the cap $e_k(\sigma_0+\sigma_1)(1+\operatorname{diam}\square)\le1$, and $\alpha_0<\frac12$ for the external-field bound of the kernel \cite[(1.16)]{Ba3} of order $m=n+n'\in\{4,5\}$, any $\alpha_0<1$ for $m\ge6$ (at the uses $m\ge\bar n+1\ge14$). The constants are independent of $\eps$, $k$, $\Omega$, $\Omega_1$, $\Omega_2$, $\tilde B$ and the torus, as \cite[(1.33)]{Ba3} asserts; they depend on $L$, which is fixed at $5$. The torus uniformity of the $O(1)$ (not in Ba{\l}aban's printed list, \cite[p.~421]{Ba3}) is the torus side of the propagator estimates of group (A): \cref{thm:a06:A5} contributes the periodic-box clause of \cite[Lemma~2.4]{Ba4}, \cref{lem:a06A:torus} gives \cite[Prop.~2.1]{Ba1} on a torus tiled by the large blocks with the constants of \cite{Ba4}, and the remaining estimates of \cref{thm:a06A:groupA} hold on the torus with the same constants, their proofs being local. With this, $[\mathrm B\text{-base}]_5$ of \cref{thm:a06:BALA}(iii) holds on every torus aligned with the blocks, with the single exception named next.

\paragraph{The walled reading.} The rotation part (\cref{sec:rot}) applies the bound \cite[(1.33)]{Ba3} once more, on tori carrying the smooth mass wall, to the two-point classes of the cut expansion of the finite order-$4$ mass counterterm and to the line-derivative families of the wall vertex: the scalar lines are $G^v_s(T_\eta,0)$ with a site potential $0\le v\le\bar v$ (in particular $v=(m^2+n\psi(t_s\cdot))t_s^2$), the photon lines are the unwalled torus propagator, the background is zero, the counterterms are the model's own, and in the two-point settings the root vertex is pinned. These classes are not per-label full classes: they are closed under the counterterm, two-vector-leg and bridgeless leg-exchange moves of \cref{def:a07:class}, but the exchange at a bridged $\phi\phi A$ subgraph carrying an external leg leaves the family. The statement needed, \cref{thm:a06:WN} below, is \cite[(1.33)]{Ba3} for such classes with walled scalar lines, with constants polynomial in $\bar v$ and free of $\eps$, $\eps_0$, the step and the torus. Its proof is the proof above run with walled lines, and it differs from the bulk at exactly two places: the zero-field kernel bounds of group (A) are replaced by their walled form (\cref{lem:a06:WN1}), and at the $\phi\phi A$ nodes the reflection identity over the full class is replaced by the clauses of \cref{thm:a06:WN}(iii), in which the exact identities of the bulk torus for the all-$C^\xi$ parts (the equality of the two leg-exchanged pinned coefficients of a bridgeless node, the vanishing of a bridged node's coefficient) are kept and the parts carrying the wall remainder $R^v$ become bounds, $|c_{S,1}(x)-c_{S,2}(x)|\le C^v_7$ for the bridgeless pair and $\le C^v_{\mathrm{BR}}$ for a bridged node paired with the counterterm member of its inner two-leg part.

% a06_WN_a: part 1 of 4 of the walled reading of Ba3 Prop. 1 (input from sec/a06_BALtorus_d7b.tex): the walled model and its zero-field kernels.
% Continues in sec/a06_WN_b.tex (classes, the two families, the statement), sec/a06_WN_c.tex (the phi phi A node clauses) and sec/a06_WN_d.tex (the bound and the constants).

\paragraph{The walled model.} Fix a step $s$, $1\le s\le K$, and work in scale-$s$ units: $\eta=L^{-s}$, $t_s\defeq L^s\eps$, $e_s\defeq e(L^s\eps)$ and $\lambda_s\defeq\lambda(L^s\eps)$ the running couplings of \cite[(1.33)]{Ba3}, and $T_\eta$ the torus, aligned with the blocks, with all sides $\ge1$. Let $v$ be a real function on the sites of $T_\eta$ with
\[
0\le v(x)\le\bar v\qquad\text{for every }x,
\]
and nothing else assumed about $v$: no smoothness, no constancy on any set, no symmetry. The walled scalar operator at zero background is
\[
H^v\defeq-\Delta^\eta_{T_\eta}+v+a_sP_s,\qquad G^v_s(T_\eta,0)\defeq(H^v)^{-1},
\]
with $a_sP_s$ the block term of the step-$s$ scalar propagator of \cite{Ba1}; at the constant $v=m^2t_s^2$ this is the bulk line $G_s(T_\eta,0)$ of Step~3(D). The wall of the rotation part is $v=(m^2+n\psi(t_s\,\cdot))t_s^2$ with a profile $0\le\psi\le1$, so that $\bar v=(m^2+n)t_s^2\le m^2+n$ once $t_s\le1$; the one-parameter family $v_\theta=(m^2+\theta n\psi(t_s\,\cdot))t_s^2$, $\theta\in[0,1]$, which interpolates between the bulk and the wall, satisfies the same bound. The \emph{walled model} is the perturbative scheme of \cite[Sect.~1]{Ba3} in which every scalar line $G_s(T_\eta,0)$, every slice \cite[(2.6)]{Ba3} and every zero-field object built from them is replaced by its $v$-analogue, while the vector lines are unchanged (the torus vector propagator carries no wall), the background is zero, and the counterterms are the model's own: a divergent two-scalar-leg subgraph $\gamma$ is renormalised by the vertex \cite[(1.7)]{Ba3} with the coefficient $\delta m^2_{s,\gamma}[v](x)$ given by Ba{\l}aban's rule \cite[p.~417]{Ba3} applied to the walled graph ($\gamma$ with both external legs at $x$ and its other leg point summed over the torus, with the same lines and the same labels, \cref{conv:a06d:LOC}, (Agr)), and the vacuum counterterm is $E_s[v]$, i.e.\ \cite[(1.24)]{Ba3} with walled lines. These counterterms depend on $x$, exactly as Ba{\l}aban's own do (\cite[p.~418]{Ba3}: the propagators ``are not translation invariant, so we have the dependence on $x$''), and the rule never uses translation invariance. Finite \cite[(1.7)]{Ba3} vertices with bounded $x$-dependent coefficients $c_v(x)$ are admitted as well; in the rotation part these are the non-renormalising parts of the inner mass counterterms and the \emph{wall vertex}, a two-scalar-leg vertex at $z$ with coefficient $-n\psi(t_sz)t_s^2$, which arises when a line of the bulk model is differentiated with respect to the wall height. The lines of a graph of the walled model are the full lines $G^v_s(T_\eta,0)$; the slices, orderings and regrouped lines of Step~3(D$_4$) are internal to the bounds.

\paragraph{Kernel notation.} For a kernel $K(x_1,x_2)$ on a lattice of spacing $\xi\in(0,1]$ and $S,T\subset\{1,2\}$, $D_S$ applies one forward difference $\partial^\xi_\mu$ in the variable $x_i$ for each $i\in S$, and $\Delta_T$ applies one H\"older quotient $|x_i-\tilde x_i|^{-\alpha}[K(\dots,x_i,\dots)-K(\dots,\tilde x_i,\dots)]$ for each $i\in T$ (in a fixed order; $\alpha\in(0,1)$ is the H\"older exponent of Step~3(D)); $\rho\defeq|x_1-x_2|$. Put $h_p(r)\defeq r^{-p}$ for $p>0$, $h_0(r)\defeq1+|\log r|$ and $h_p(r)\defeq1$ for $p<0$. In this notation the bounds of \cref{lem:a06d:S47}(b) on a remainder kernel $R$ read
\begin{equation}\label{eq:a06:WNR}
|\Delta_TD_SR(x_1,x_2)|\le C\,h_{|S|+\alpha|T|-1}(\xi+\rho)\,e^{-\delta\rho}\prod_{i\in T}|x_i-\tilde x_i|^\alpha,\qquad|S|,|T|\le2,
\end{equation}
and the zero-field, zero-mass form of the group-(A) estimates (\cite[Lemma~2.4, (2.35)--(2.36)]{Ba4}, \cref{thm:a06:A5}), extended to a difference in each variable by the same contour-shift argument, one factor $(e^{i\xi p_\mu}-1)/\xi$ per difference, for the pinned free operator $T\defeq-\Delta^\xi_X+bP_Y$ ($X$ a union of unit blocks with Neumann conditions, a periodic box or $\xi\Z^3$; $Y$ the unit blocks, or the $L$-blocks when $X$ is a union of $L$-blocks; $P_Y$ the block-average projection; $b\in[b_-,b_+]\subset(0,\infty)$) reads
\begin{equation}\label{eq:a06:WNZ1}
|\Delta_TD_ST^{-1}(x_1,x_2)|\le C_1(\xi+\rho)^{-p(S,T)}e^{-\delta\rho}\prod_{i\in T}|x_i-\tilde x_i|^\alpha,\qquad p(S,T)\defeq1+|S|+\alpha|T|,
\end{equation}
at every pair of points of $X$, with $T\ge\gamma'>0$ and $C_1,\delta,\gamma'$ depending only on $d=3$, $L$, $b_\pm$ and $\alpha$. (A quotient over two points at distance $>1$ is bounded by the two values, so \eqref{eq:a06:WNZ1} is needed only at $|x_i-\tilde x_i|\le1$; the same convention applies to every bound below.)

\begin{lemma}[lattice convolution]\label{lem:a06:WNK}
Let $0\le p_1,p_2<3$ and $\delta>0$. Then for all $x,y\in\xi\Z^3$ and $\xi\in(0,1]$,
\[
\sum_{z\in\xi\Z^3}\xi^3(\xi+|x-z|)^{-p_1}e^{-\delta|x-z|}(\xi+|z-y|)^{-p_2}e^{-\delta|z-y|}\ \le\ C_K\,h_{p_1+p_2-3}(\xi+|x-y|)\,e^{-\delta|x-y|/2},
\]
with $C_K=C_K(p_1,p_2,\delta)$. The same holds on a torus or a Neumann domain, the sum running over its sites.
\end{lemma}

\begin{proof}
Put $r\defeq|x-y|$. Since $|x-z|+|z-y|\ge r$, the exponentials are $\le e^{-\delta r/2}e^{-\delta(|x-z|+|z-y|)/2}$. On $\{|x-z|\le r/2\}$ one has $\xi+|z-y|\ge(\xi+r)/2$, and $\sum_{|x-z|\le r/2}\xi^3(\xi+|x-z|)^{-p_1}\le C(\xi+r)^{3-p_1}$ because $p_1<3$; the region $\{|z-y|\le r/2\}$ is symmetric; on the rest both factors are $\le C(\xi+|z-x|)^{-(p_1+p_2)}$ up to the exponentials, and $\sum_{|z|\ge r/2}\xi^3(\xi+|z|)^{-(p_1+p_2)}e^{-\delta|z|/2}$ is $\le C(\xi+r)^{3-p_1-p_2}$ if $p_1+p_2>3$, $\le C(1+|\log(\xi+r)|)$ if $p_1+p_2=3$ and $\le C$ otherwise. For $\xi+r\ge1$ the exponential factor dominates. On a torus or a Neumann domain the sites are a subset of $\xi\Z^3$ with the quotient (resp.\ restricted) distance, and the same sums bound the result.
\end{proof}

\begin{lemma}[walled zero-field kernels]\label{lem:a06:WN1}
Let $X$, $Y$, $b$ and $T=-\Delta^\xi_X+bP_Y$ be as in \eqref{eq:a06:WNZ1}, let $w\colon X\to[0,\bar v]$ and $T_w\defeq T+w$. Then:
\begin{enumerate}[label=\textup{(\alph*)},nosep,leftmargin=*]
\item\textup{(algebra)} $T_w\ge T\ge\gamma'$. The function $w$ commutes with every multiplication operator ($e^{\beta\rho}$ for a Lipschitz $\rho$, cutoffs, the gauge factors $V_\lambda$) and cancels in every difference of two operators carrying the same $w$; hence the gauge covariance $H^w(A^\lambda)=V_\lambda H^w(A)V_\lambda^*$, the commutator identities with cutoffs, the slice identity \cite[(2.6)]{Ba3} and its rescaling hold for the walled operator verbatim, the rescaled potential at scale $s_j$ being $s_j^2w(s_j\,\cdot)\in[0,\bar v s_j^2]\subset[0,\bar v]$, and the slices of the walled propagator are $G^{v,\eta}_{(j)}=a_j^2s_j^{-4}G^v_jQ_j^*C^{(j),v}Q_jG^v_j$ with $G^v_j\defeq(-\Delta+v+a_js_j^{-2}P_j)^{-1}$ and $C^{(j),v}$ the block operator built from $G^v_j$. Every walled zero-field scalar kernel is a scalar kernel times $\mathbf 1_N$, because $w$ is a scalar.
\item\textup{(the walled kernel bound)} $T_w^{-1}$ satisfies \eqref{eq:a06:WNZ1} at every pair of points, with $C_1$ replaced by $C_1^w\defeq C(d,L,\alpha,b_+,\gamma',\delta,C_1)(1+\bar v)^2$ and $\delta$ replaced by $\delta^w\defeq\min\{\delta/4,\beta'/\sqrt3\}/4$, where $\beta'\le1$ is the Combes--Thomas rate of \cref{lem:a07:CT} for $T_w$ (it depends only on $\gamma'$, $b_+$ and $L$, not on $w$).
\item\textup{(block sums)} For $|S|,|T|\le1$ and every block $B(y)$: $|\Delta_TD_S(T_w^{-1}\mathbf 1_{B(y)})(x)|\le c^we^{-\delta^w|x-y|}\prod_{i\in T}|x-\tilde x|^\alpha$.
\item\textup{(the remainder)} On $\xi\Z^3$ or a periodic box let $C^\xi\defeq(-\Delta^\xi+1)^{-1}$ be the free unit-mass kernel of Step~3(D$_4$) and $R^w\defeq T_w^{-1}-C^\xi$. Then $R^w$ satisfies \eqref{eq:a06:WNR} at every pair of points with $C=C_R(1+\bar v)^3$ and $\delta=\delta_R$, $C_R$ and $\delta_R$ depending only on $d,L,\alpha,b_\pm,\gamma'$ and the constants of \eqref{eq:a06:WNZ1}. In particular $R^w$ is a remainder line in the sense of \cref{lem:a06d:S47}(b), with constant $C_R(1+\bar v)^3$, and so are its slices and the kernels of \cref{lem:a06d:S47}(c) built from it.
\item\textup{(slices, cube kernels)} The slice bounds \cite[(2.10)--(2.12)]{Ba3} and the kernel bounds of the cube step of Step~3(D) hold for the walled slices $G^{v,\eta}_{(j)}$, including the two-ended entries and the mixed second difference, with constants depending polynomially on $\bar v$ and otherwise as in the bulk; the Combes--Thomas bound of \cref{lem:a07:CT} holds for $H^v$ with the $v$-free constants.
\item\textup{(periodisation)} $G^v_s(T_\eta,0)(x,x')=\sum_{n}G^{v_{\mathrm{per}}}_s(\eta\Z^3,0)(\tilde x,\tilde x'+n)$, the sum over the period lattice of $T_\eta$, $\tilde x,\tilde x'$ lifts of $x,x'$ and $v_{\mathrm{per}}$ the periodic extension of $v$; the sum converges absolutely, and the same identity holds slice by slice.
\end{enumerate}
\end{lemma}

\begin{proof}
(a) $w$ is a real scalar multiplication operator: it is self-adjoint and $\ge0$, it commutes with every other multiplication operator, and it cancels in $H^w(A)-H^w(A')$ and in $[H^w,\chi]$. The slice identity is the resolvent algebra of $-\Delta+w$ with the pin terms, which uses no property of the zeroth-order operator beyond coercivity; rescaling by $s_j$ multiplies a potential by $s_j^2$ and composes it with the dilation. Since $v$ acts as $v\mathbf 1_N$, the inverse of $H^v$ at zero field is a scalar kernel times $\mathbf 1_N$.

(b) \emph{Step 1 (Combes--Thomas).} $T_w=D^*D+w+bP_Y$ is of the form of \cref{lem:a07:CT} with $V=w\ge0$, $\Pi_B$ the block-average projections and $T_w\ge\gamma'$; so $\norm{e^{\beta\rho}T_w^{-1}e^{-\beta\rho}}\le2/\gamma'$ for every $1$-Lipschitz $\rho$ and $0\le\beta\le\beta'$, $\rho$ being $\sqrt3$-Lipschitz for the sup-metric, so that $\beta\le\beta'/\sqrt3$ is imposed, with $\beta'$ depending only on $\gamma'$, $b_+$ and the block size.
\emph{Step 2 (identity).} $T_w^{-1}=T^{-1}-T^{-1}wT^{-1}+T^{-1}wT_w^{-1}wT^{-1}$, which combines $T_w^{-1}=T^{-1}-T^{-1}wT_w^{-1}$ with $T_w^{-1}=T^{-1}-T_w^{-1}wT^{-1}$.
\emph{Step 3 (no operations).} Put $r\defeq|x-y|$. The first term is $\le C_1(\xi+r)^{-1}e^{-\delta r}$ by \eqref{eq:a06:WNZ1}. The second is $\le\bar vC_1^2C_Kh_{-1}(\xi+r)e^{-\delta r/2}=\bar vC_1^2C_Ke^{-\delta r/2}$ by \cref{lem:a06:WNK} with $p_1=p_2=1$. For the third write $a(z)\defeq T^{-1}(x,z)w(z)$, $b(z')\defeq w(z')T^{-1}(z',y)$ and $\rho\defeq|\cdot-x|$; then $|\ip a{T_w^{-1}b}|\le\norm{e^{\beta\rho}a}_2\,(2/\gamma')\,\norm{e^{-\beta\rho}b}_2$, and since $|z'-x|\ge r-|z'-y|$, $\norm{e^{-\beta\rho}b}_2\le e^{-\beta r}\norm{e^{\beta|\cdot-y|}b}_2$; both weighted norms are $\le\bar vC_1J$ with $J^2\defeq\sup_\xi\sum_z\xi^3(\xi+|z|)^{-2}e^{-(\delta-2\beta)|z|}<\infty$ (here $d=3$ and $\beta\le\delta/4$). Hence $|T_w^{-1}(x,y)|\le C^{(1)}(1+\bar v)^2(\xi+r)^{-1}e^{-\delta'r}$ with $\delta'\defeq\min\{\delta/2,\beta\}$ (for $\xi+r\le1$ the constant terms are $\le(\xi+r)^{-1}$, for $\xi+r\ge1$ the exponential absorbs the power).
\emph{Step 4 (general $S,T$).} Apply $\Delta_TD_S$ to the identity of step 2. The operations at $x_1$ act only on the first factor $T^{-1}(x_1,\cdot)$ and those at $x_2$ only on the last factor $T^{-1}(\cdot,x_2)$; $w$ and the middle factor $T_w^{-1}$ are never differentiated. Put $p_i\defeq1+[i\in S]+\alpha[i\in T]\in[1,2+\alpha]$, so that $p(S,T)=p_1+p_2-1$. The first term is \eqref{eq:a06:WNZ1}. The second is, by \eqref{eq:a06:WNZ1} for the two factors and \cref{lem:a06:WNK}, $\le2^{|T|}\bar vC_1^2C_Kh_{p_1+p_2-3}(\xi+\rho)e^{-\delta\rho/2}\prod|x_i-\tilde x_i|^\alpha$. For the third insert the bound of step 3 for the middle factor and apply \cref{lem:a06:WNK} twice: the convolution of $(\xi+|x_1-z|)^{-p_1}$ with $(\xi+|z-z'|)^{-1}$ gives $h_{p_1-2}(\xi+|x_1-z'|)$, which is $\le C_\alpha(\xi+|x_1-z'|)^{-q}$ with $q\defeq\max\{p_1-2,0\}+\alpha/2<3$ (this covers the logarithm at $p_1=2$), and the second convolution with $(\xi+|z'-x_2|)^{-p_2}$ gives $h_{q+p_2-3}(\xi+\rho)$. Since $p_1+p_2-3<p(S,T)$ and $q+p_2-3<p(S,T)$ with $p(S,T)\ge1$, both $h$-factors are $\le C(\xi+\rho)^{-p(S,T)}$ for $\xi+\rho\le1$, and for $\xi+\rho\ge1$ the exponentials dominate. Each application of \cref{lem:a06:WNK} halves the rate: the second term decays at the rate $\delta/2$ and the third at $\delta'/4$. Taking $\beta\defeq\min\{\delta/4,\beta'/\sqrt3\}$, so that $\delta'=\beta$, and collecting the constants gives $C_1^w=C(1+\bar v)^2$ and the rate $\delta^w=\min\{\delta/2,\beta/4\}=\beta/4$.

(c) Sum (b) over $x'\in B(y)$: since $p(S,T)\le2+\alpha<3$, $\sum_{x'\in B(y)}\xi^3(\xi+|x-x'|)^{-p}e^{-\delta^w|x-x'|}\le C(p)e^{-\delta^w\dist(x,B(y))}$.

(d) Since $(C^\xi)^{-1}=-\Delta^\xi+1$ and $T_w^{-1}$ exists, $R^w=T_w^{-1}\big((C^\xi)^{-1}-T_w\big)C^\xi=T_w^{-1}(1-w)C^\xi-b\,T_w^{-1}P_YC^\xi$. This is the identity \cite[p.~437]{Ba3} ($G^\xi(0)=C^\xi+G^\xi(0)(1-m^2-aP)C^\xi$) with the constant $1-m^2$ replaced by the function $1-w(z)$, which is never differentiated: the operations at $x_1$ fall on $T_w^{-1}(x_1,\cdot)$, those at $x_2$ on $C^\xi(\cdot,x_2)$. In the first term, (b) for $T_w^{-1}$, the free bound \eqref{eq:a06:WNZ1} for $C^\xi$ (it is the case $b=0$, $m^2=1$ of the group-(A) estimates), $|1-w|\le1+\bar v$ and \cref{lem:a06:WNK} give $C(1+\bar v)^3h_{p_1+p_2-3}(\xi+\rho)e^{-\delta_R\rho}\prod|x_i-\tilde x_i|^\alpha$ with $p_1+p_2-3=|S|+\alpha|T|-1$. The second term is $b\sum_Y|Y|^{-1}(T_w^{-1}\mathbf 1_Y)(x_1)(\mathbf 1_Y^{\mathsf T}C^\xi)(x_2)$, and (c) for $T_w^{-1}\mathbf 1_Y$ together with the block-sum bound for $C^\xi\mathbf 1_Y$ give $Ce^{-\delta_R\rho}\prod|x_i-\tilde x_i|^\alpha$ after the sum over $Y$, which is $\le h_{|S|+\alpha|T|-1}(\xi+\rho)$ for $\xi+\rho\le1$. The slices of $R^w$ and the kernels of \cref{lem:a06d:S47}(c) are obtained from these bounds exactly as in \cref{lem:a06d:S47}, which uses only the bounds \eqref{eq:a06:WNR}.

(e) The slice $G^{v,\eta}_{(j)}$ is the composition of (a) with $G^v_j$ bounded by (b) at scale $j$ (unit-block units after rescaling, with the potential $s_j^2v(s_j\,\cdot)\in[0,\bar v]$) and with the block-operator kernel $C^{(j),v}$, whose exponential bound uses only the coercivity $H^v\ge H^0\ge\gamma_0$ and the Combes--Thomas argument of \cref{lem:a07:CT}, so its constants are $v$-free; the bulk proof of \cite[(2.10)--(2.12)]{Ba3} in Step~3(D) is a composition of these three bounds and of lattice sums, and the mixed second difference and the two-ended entries are the cases $|S|=2$ of (b). The kernel bounds of the cube step are sums of slice bounds over the scales. The Combes--Thomas bound is \cref{lem:a07:CT} with $V=v$.

(f) The right side is well defined because the image sum converges absolutely by (b) on $\eta\Z^3$ with the periodic potential $v_{\mathrm{per}}$ (which lies in $[0,\bar v]$); it is periodic, and it solves the torus equation because the $\eta\Z^3$ operator with $v_{\mathrm{per}}$ commutes with the period translations; the torus inverse is unique. The slice identity (a) is compatible with the periodisation term by term.
\end{proof}

% a06_WN_b: part 2 of 4 of the walled reading of Ba3 Prop. 1 (input from sec/a06_BALtorus_d7b.tex): classes closed under the moves, the two families of the rotation part, and the statement.

\paragraph{Classes closed under the moves.} Recall the moves (P), (P$'$), (W), (X) and the classes of \cref{def:a07:class}, and the cut terms and orbits of \cref{def:a07:cut}. A \emph{decorated graph} of the walled model is a graph of the model (lines $G^v_s(T_\eta,0)$ and the torus vector propagator, vertices \cite[(1.6)--(1.15)]{Ba3} with their labels (Loc), counterterm vertices with the model's own coefficients, finite \cite[(1.7)]{Ba3} vertices with bounded coefficients) whose external legs carry external fields which may be kernels of several variables \cite[p.~419]{Ba3}, in particular the two ends of a cut line; these legs are the \emph{slots} of the graph. A finite set $\cC$ of decorated graphs with a common slot structure (the same slots, with the same kinds, differentiations, labels and fields), a common order, a common label multiset and a common multiset of finite-vertex coefficients is \emph{closed under the moves} if every move of the set $M_0$, consisting of (P) and (P$'$) at a divergent two-scalar-leg subgraph of a member, (W) at a divergent two-vector-leg subgraph, and (X) at a \emph{bridgeless} divergent $\phi\phi A$ subgraph (at the member itself included, when the member is of one of these types), gives a member. The set $M_0$ is a subset of the moves of \cref{def:a07:class} and has the properties (A1)--(A5) of that definition, which are properties of the individual moves; so \cref{lem:a07:orb} applies with $M=M_0$. Per-label full classes are closed under the moves (\cref{def:a07:class}), and so is every class $\Cl_{M}(K)$, $M\supset M_0$, of a connected component $K$ of a cut graph in \cref{lem:a07:orb}(c), with the cut ends as slots. The restriction to bridgeless subgraphs in (X) is forced by the two-point settings: the exchange at a bridged divergent $\phi\phi A$ subgraph carrying an external leg produces a one-particle-reducible graph, which is not a term of a family of \cite[(1.23)]{Ba3} graphs (\cref{thm:a06:WN}(iii)(c) shows that no partner is needed there). In a \emph{two-point setting} the members have two external scalar legs, undifferentiated, one of them, the \emph{root}, at a vertex pinned at the point $x$, the other at a vertex that is summed; the coefficient is a function of $x$, and the bounds below hold uniformly in $x$. A $\phi\phi A$ subgraph $S$ is \emph{bridgeless} if no internal line of $S$ disconnects it (a self-loop is not a bridge), and \emph{bridged} otherwise.

\paragraph{The two families of the rotation part.} (a) \emph{The cut expansion of the finite order-$4$ mass counterterm.} The order-$4$ two-point coefficient of the walled model is computed once with the $\eps$-level lines $C^v_0$, $C^\gamma$ and once with the step-$K$ lines $G^v_K$, $G^\gamma_K$; the difference, with the scalar lines decomposed as $C^v_0=G^v_K+R^v_K$ and the vector lines as $C^\gamma=G^\gamma_K+R^\gamma_K$, is the sum \eqref{eq:a07:tel} of \cref{def:a07:cut} over the cut terms with at least one cut line, the cut ends carrying the kernels $R^v_K$, $R^\gamma_K$. By \cref{lem:a07:orb}(c),(d) each orbit is a product of classes $\Cl_M(K)$ of the connected components $K$ of its cut graph, each closed under the moves, with the cut ends as slots; the component carrying the root is a two-point class with the root pinned at $x$, the others are vacuum components with slots. (b) \emph{The line-derivative family of the wall vertex.} Let $\cF^{(4)}$ be the set of two-point graphs of \cite[(1.23)]{Ba3} of coupling order $4$: one-particle-irreducible, two undifferentiated external scalar legs, the root at $x$, the inner \cite[(1.7)]{Ba3} vertices carrying the model's own step-$s$ counterterm of their order and their non-renormalising parts as finite vertices. For $\Gamma\in\cF^{(4)}$ and a scalar line $l$ of $\Gamma$ let $\Gamma_{l\to v}$ be $\Gamma$ with $l$ subdivided by a two-scalar-leg vertex $v$ at a point $z$ with a bounded coefficient $c(z)$ (the wall vertex), and
\[
(\Gamma_{l\to v})_{\mathrm{ren}}\ \defeq\ \sum_{F}\prod_{\gamma\in F}(-T_\gamma)\,\Gamma_{l\to v},
\]
the sum over the forests $F$ of divergent two-scalar-leg subgraphs of $\Gamma$ avoiding $l$ (i.e.\ not containing $l$), $T_\gamma$ the counterterm operation of (Agr). The family is $\sum_{\Gamma\in\cF^{(4)}}\sum_l(\Gamma_{l\to v})_{\mathrm{ren}}$, with no cut line.

\begin{lemma}[the line-derivative family]\label{lem:a06:WNNC}
\textup{(a)} Every connected subgraph of $\Gamma_{l\to v}$ that contains $v$ has degree $D\ge1$ or vanishes. Hence the divergent subgraphs of $\Gamma_{l\to v}$ are exactly the divergent subgraphs of $\Gamma$ avoiding $l$, and $(\Gamma_{l\to v})_{\mathrm{ren}}$ is the renormalised forest sum of $\Gamma_{l\to v}$. \textup{(b)} A forest of divergent two-scalar-leg subgraphs of $\Gamma\in\cF^{(4)}$ avoiding $l$ has at most one element. \textup{(c)} Group the terms $T_F\Gamma_{l\to v}\defeq\prod_{\gamma\in F}(-T_\gamma)\Gamma_{l\to v}$ ($\Gamma\in\cF^{(4)}$, $l$ a scalar line of $\Gamma$, $F$ as above) by label multiset and by the multiset of finite-vertex coefficients, the wall vertex included. Each group is a two-point class closed under the moves, with the root pinned at $x$.
\end{lemma}

\begin{proof}
(a) The vertex $v$ is a \cite[(1.7)]{Ba3} vertex with a bounded coefficient, so it has $D(v)=2$ and every connected subgraph $H$ containing it has $D(H)\ge1$ if $H$ has at least two slots (\cref{lem:a06d:S25}(b), \cite[Cor.~2.3]{Ba3}). A connected subgraph with exactly one slot vanishes: with a scalar slot by the symmetry $\phi\to-\phi$ of every vertex (its amputated kernel $K_a$ with one open scalar index satisfies $K_a=-K_a$), with a vector slot at zero background because its scalar index loops close inside it and $\Tr q^{\mathrm{odd}}=0$ (Step~3(D$_4$)); and a connected subgraph of $\Gamma_{l\to v}$ with no slot is $\Gamma_{l\to v}$ itself, which has two. The divergent subgraphs of $\Gamma_{l\to v}$ therefore avoid $v$, i.e.\ are subgraphs of $\Gamma$ not containing $l$, with the same vertices, lines, slots and degree (\cref{lem:a07:orb}(a)); conversely a divergent subgraph of $\Gamma$ avoiding $l$ is a subgraph of $\Gamma_{l\to v}$ with the same data.

(b) A proper divergent two-scalar-leg subgraph of $\Gamma$ has coupling order $2$ (\cite[(1.23)]{Ba3}; the order-$4$ one is $\Gamma$ itself, which contains $l$). Two such subgraphs cannot be nested, and two disjoint ones $\gamma,\gamma'$ would exhaust the vertices of $\Gamma$: each has exactly two slots, the two external legs of $\Gamma$ are among the four slots, so at most one line joins $\gamma$ to $\gamma'$; if exactly one does it is a bridge of $\Gamma$, and if none does $\Gamma$ is disconnected; both contradict $\Gamma\in\cF^{(4)}$.

(c) Fix a group. \emph{(P), (P$'$).} A divergent two-scalar-leg subgraph $\gamma$ of a member $T_F\Gamma_{l\to v}$ avoids $v$ and $l$ by (a); the move replaces the term with forest $F$ by the term with forest $F\cup\{\gamma\}$ (or $F\setminus\{\gamma\}$), a forest avoiding $l$ by (b), of the same $\Gamma$ and $l$; the counterterm vertex carries $\gamma$'s labels (Agr) and the model's own coefficient, so the labels and the finite-vertex coefficients of the group are kept. \emph{(W).} A divergent two-vector-leg subgraph $H$ of a member avoids $v$ and $l$ by (a). Replacing $H$ by a graph $H'$ of its per-label two-vector-leg class (same two vector slots, same order, same label multiset, its counterterm members included) changes nothing outside $H$: the line $l$ with the vertex $v$ on it and the two external scalar legs are untouched, the order and the labels are kept, and the result is one-particle-irreducible: since $T_F\Gamma_{l\to v}$ is one-particle-irreducible, the two attachment points of $H$ are connected outside $H$ (else $H$'s vector line to the rest would be a bridge), so an internal bridge of $H'$ separating its two slots is not a bridge of the new graph, and an internal bridge of $H'$ with a slot-free side bounds a one-slot subgraph, which vanishes by (a). Smoothing $v$ back into a line gives a graph of $\cF^{(4)}$ with the counterterm vertices of $F$ inside or outside $H$ as the case may be (a forest element inside $H$ is a counterterm member of $H$'s class), so the result is a member of the group. \emph{(X).} The exchange at a bridgeless $\phi\phi A$ subgraph $P$ moves only the attachments at $P$'s two scalar slots and no vertex: the line $l$ keeps its ends up to this re-attachment and $v$ stays on it, the external legs stay undifferentiated (both slots of $P$ carry no difference, since $P$ is divergent), and the result is one-particle-irreducible: for a line $\ell\notin P$, $P$ is connected in both graphs minus $\ell$ with the same multiset of outer attachment points, and for $\ell\in P$, $P\setminus\ell$ is connected because $P$ is bridgeless. The forest element, if any, is nested with $P$ or disjoint from it (a divergent two-scalar-leg subgraph crossing a bridgeless order-$3$ $\phi\phi A$ subgraph would meet it in a single one-vector-leg vertex whose two internal lines of $P$ would be its two scalar slots, leaving its vector line as a bridge of the forest element); a forest element sharing a vertex with $P$ but no line (the tadpole at the \cite[(1.6)]{Ba3} vertex of the double-line shape) is untouched, since \textup{(X)} moves attachments only; so $F$ is carried along. \emph{Bridged $\phi\phi A$ subgraphs.} No move of $M_0$ acts there. (The exchange at a bridged divergent $\phi\phi A$ subgraph $\{u\}\cup S_2$ of a member, with the external leg of the member at $u$, would move that leg to $S_2$, which then hangs on the bridge alone: the result is one-particle-reducible and is not a term of the family. Such a member exists: $u$, $u'$ one-vector-leg vertices \cite[(1.8)]{Ba3} joined by a vector line, $w$ a \cite[(1.6)]{Ba3} vertex with a self-loop, scalar lines $u$--$w$ and $w$--$u'$, the root at $u$ and the wall vertex on $w$--$u'$; the node $\{u,w\}$ is a dangerous bridged $\phi\phi A$ node of degree $0$.) Finally every member has its root pinned at $x$, since no move moves a vertex.
\end{proof}

\begin{theorem}[the walled form of {\cite[Prop.~1]{Ba3}} for classes closed under the moves]\label{thm:a06:WN}
In the walled model at a step $s\le K$ (potential $0\le v\le\bar v$ on the scalar lines only, vector lines unwalled, zero background, the model's own counterterms, finite \cite[(1.7)]{Ba3} vertices with bounded coefficients, torus aligned with the blocks with sides $\ge1$) let $\cC$ be a class closed under the moves, of order $\le\bar n$, with slot fields $\Phi_{\mathrm{ext}}$, $A_{\mathrm{ext}}$ and finite-vertex coefficients $c_v$. Then:
\begin{enumerate}[label=\textup{(\roman*)},nosep,leftmargin=*]
\item\textup{(regrouping)} \Cref{lem:a07:orb} (with $M=M_0$ in the two-point settings, under which the families of \cite[(1.23)]{Ba3} graphs are closed) and \cref{thm:a06d:S23} hold verbatim with walled lines; in particular the cut expansion of the finite order-$4$ mass counterterm is a finite sum over orbits, each orbit the contraction through the kernels $R^v_K$, $R^\gamma_K$ of per-component classes closed under the moves, with the invariants of \cref{lem:a07:orb}(b),(f); and the line-derivative family of the wall vertex is a finite disjoint union of two-point classes closed under the moves (\cref{lem:a06:WNNC}).
\item\textup{(the bound)} With $e_s=e(L^s\eps)$ and $\lambda_s=\lambda(L^s\eps)$,
\[
\big|E^v(\cC;\{\square(v)\},\Phi_{\mathrm{ext}},A_{\mathrm{ext}})\big|\ \le\ O(1)^v\,e_s^{\,d_v}\lambda_s^{\,d_s}\prod_{v\ \mathrm{finite}}\sup_z|c_v(z)|\;e^{-\delta^v_0d(\{\square(v)\})}\prod\norm{h\Phi_{\mathrm{ext}}}_{1,\alpha_0}\prod\norm{h'A_{\mathrm{ext}}}_{1,\alpha_0},
\]
the form \cite[(1.33)]{Ba3} ($E^v$ the class sum of \cref{lem:a07:orb}(d) with walled lines): $d_v$, $d_s$ count the couplings of the action vertices, a counterterm vertex counting as the exact counterterm of the order it replaces; the localisations $h$, $h'$ are those of (H$_A$); the slot fields are kernels of several variables, each measured by its norm \cite[(1.32)]{Ba3} in the variable at the slot; in a two-point setting the left side is the coefficient at the pinned root $x$, and the bound is uniform in $x$. The constants $O(1)^v$ and $\delta^v_0$ depend only on the order, on the parameters of Step~3(D) ($d=3$, $L=5$, $a$, $N$, $\norm q$, $\bar n$, $\alpha_0$, $m$, $\lmu$, $M$, $R$) and polynomially on $\bar v$; they do not depend on $\eps$, on $\eps_0\in(0,\eps_0^*]$ below the caps of Step~3, on $s$, $K$, the labels, the torus side, $x$, or on $v$ beyond $\bar v$.
\item\textup{(the $\phi\phi A$ node clauses)} Let $S$ be a dangerous $\phi\phi A$ node of a member $G$ of $\cC$ along an ordering, and let $c$ be its regrouped zeroth-order local coefficient in the form of \cref{thm:a06d:S48}: rescaled to $\xi\Z^3$, lines $\breve G_l=C^\xi+R_l$ with $R_l=R^v_l$ of \cref{lem:a06:WN1}(d) on the scalar lines and the bulk remainder on the vector lines, multipliers at the anchor, and the decomposition $c=c^C+c^R$ into the all-$C^\xi$ part and the terms with at least one $R$-line.
\begin{enumerate}[label=\textup{(\alph*)},nosep,leftmargin=*]
\item $S$ bridgeless, $G^{\mathrm{sw}}$ the member obtained by \textup{(X)} at $S$, a common anchor in $S$: $c(S\text{ in }G)+c(S\text{ in }G^{\mathrm{sw}})=0$ exactly, for any line kernels.
\item $S$ bridgeless with its two scalar slots at distinct vertices $u_1\ne u_2$, the pinned root at $u_1$ in $G$ and at $u_2$ in $G^{\mathrm{sw}}$, anchors at the root; $c_{S,1}(x)$, $c_{S,2}(x)$ the two coefficients: $c^C_{S,1}(x)=c^C_{S,2}(x)$, and $|c_{S,1}(x)-c_{S,2}(x)|\le C^v_7$.
\item $S$ bridged. Then $S=\{u\}\cup S_2$, where $u$ is a one-vector-leg vertex \cite[(1.8)]{Ba3} carrying one scalar slot and the vector slot of $S$, its difference acting on the bridge $l_0=(u,w)$, $w\in S_2$, and $S_2$ is a two-scalar-leg subgraph of order $2$ (the exchange bubble, the \cite[(1.6)]{Ba3} tadpole, a seagull with a vector loop, or a \cite[(1.7)]{Ba3} vertex); these are the shapes \cite[(2.21b)--(g)]{Ba3}. If $u$ is summed, then with the anchor in $S_2$, $c=0$ for every member and any line kernels. If $u$ carries the pinned root (anchor $u$ at $x$), then $c^C=0$ for every member, and the $R$-part of the pair $(S\text{ with }S_2)+(S\text{ with }T_{S_2})$, both members of $\cC$ by \textup{(P)}, is $\le C^v_{\mathrm{BR}}$ in absolute value; the $R$-part of a single member is in general not bounded uniformly in $\xi$.
\end{enumerate}
$C^v_7$ and $C^v_{\mathrm{BR}}$ are the $\phi\phi A$ constant of \cref{thm:a06d:S48} times a polynomial in $\bar v$, free of $\xi$, $j_0$, $s$, $\eps$, the labels, the torus and $x$.
\end{enumerate}
\end{theorem}

The rotation part consumes (ii) and (iii) at the step $s=K$ for the per-component classes of the cut expansion in (i), and at every step $s=a\le K$ for the line-derivative family, both with $v$ in the family $v_\theta$ and at coupling order $\le4$. The proof of (ii) is by induction on the number of vertices of the members, as the proof of \cite[(1.33)]{Ba3} is (\cref{lem:a06d:S25}); the induction runs over classes closed under the moves, the completed pieces of a cut partnered subgraph being such classes of lower order.

\begin{proof}[Proof of \textup{(i)}]
\Cref{lem:a07:orb} uses the lines only through the properties (A1)--(A5) of the moves, the multilinearity of the expansion \eqref{eq:a07:tel} and the symmetry of the decomposed kernels; $R^v_K=C^v_0-G^v_K$ is a difference of two symmetric resolvents, hence symmetric. In the two-point setting the family of \cite[(1.23)]{Ba3} graphs is closed under $M_0$: (P) contracts or expands a one-particle-irreducible two-leg subgraph (a bridge of the new graph would be one of the old), and (W), (X) at bridgeless subgraphs keep one-particle-irreducibility and the undifferentiated external legs by the arguments given in the proof of \cref{lem:a06:WNNC}(c). \Cref{thm:a06d:S23} and \cref{lem:a06d:S22} are statements about labels and the cut rule. Degrees count lines (CV). So (i) is the bulk statement, read with walled lines, and the second half is \cref{lem:a06:WNNC}(c).
\end{proof}

% a06_WN_c: part 3 of 4 of the walled reading of Ba3 Prop. 1 (input from sec/a06_BALtorus_d7b.tex): the proof of thm:a06:WN(iii), the phi phi A node clauses with walled lines.

\begin{proof}[Proof of \cref{thm:a06:WN}\textup{(iii)}]
\emph{The shapes.} A dangerous $\phi\phi A$ node has order $3$ and one loop (Step~3(D$_4$)), and a node containing a finite \cite[(1.7)]{Ba3} vertex is not dangerous (\cref{lem:a06d:S25}(b)); by \cite[p.~429]{Ba3} its vertices are of the types \cite[(1.6)--(1.8), (1.10)]{Ba3} or divergent-order counterterm vertices. With an odd number of vector legs the vertex multisets are three one-vector-leg vertices \cite[(1.8)]{Ba3}; one one-vector-leg vertex with one two-vector-leg vertex of type \cite[(1.8)]{Ba3} or \cite[(1.10)]{Ba3}; one three-vector-leg vertex \cite[(1.8)]{Ba3}; or one one-vector-leg vertex with a \cite[(1.6)]{Ba3} vertex or with an order-$2$ counterterm vertex. The line count gives one loop in each case, so the internal lines form one cycle through all vertices (a triangle or a double line) in the bridgeless case, and a bridge otherwise. The single three-vector-leg vertex is excluded (a vector self-loop leaves a differentiated scalar slot, a scalar self-loop leaves three vector slots). The bridgeless shapes are: \emph{T}, three one-vector-leg vertices in a triangle, the apex carrying the vector slot and the two base vertices the scalar slots, which are their undifferentiated legs, the base vertices' differences acting on their lines to the apex, and the apex's difference on one of its two lines; \emph{B8}, resp.\ \emph{B10}, a one-vector-leg vertex $u$ and a two-vector-leg vertex $w$ of type \cite[(1.8)]{Ba3}, resp.\ \cite[(1.10)]{Ba3}, joined by a vector line and a scalar line, the slots being $u$'s undifferentiated scalar leg, $w$'s undifferentiated scalar leg and $w$'s second vector leg; \emph{BL}, a one-vector-leg vertex $u$ and a \cite[(1.6)]{Ba3} vertex $w$ joined by two scalar lines ending at legs of different index pairs of $w$, the slots being the two remaining legs of $w$ and $u$'s vector leg. In each shape every scalar slot is an undifferentiated leg and the vector slot carries no difference, as the degree $D=0$ requires. The bridged shapes are \cite[(2.21a)--(g)]{Ba3}; \cite[(2.21a)]{Ba3}, a vector bridge whose other side is a one-vector-leg vertex with a scalar loop, vanishes identically, and in the remaining shapes: a side of the bridge with no slot would be a one-slot subgraph, which vanishes (\cref{lem:a06:WNNC}(a)); a vector bridge is excluded: the side carrying the vector slot has two vector legs (the slot and the bridge), hence order $\ge2$, so the other side has order $1$ and is a single one-vector-leg vertex with no loop, whose two scalar legs are then both external, i.e.\ both scalar slots lie on it, and its difference acts on a slot, against $D=0$; so the bridge $l_0$ is scalar, each side carries one scalar slot, the side with the vector slot has order $\ge1$ and the other, a two-scalar-leg subgraph, order $\ge2$, whence orders $1+2$: the first side is a single one-vector-leg vertex $u$ whose difference, not allowed on its scalar slot or its vector leg, acts on $l_0$, and $S_2$ is one of the order-$2$ two-scalar-leg graphs, none of which has a bridge. This is the structure asserted in (c), and it is what Ba{\l}aban's arrows in \cite[(2.21b)--(g)]{Ba3} show.

\emph{(a): the common-anchor cancellation.} Let $\sigma_i\in u_i$ be the two scalar slots of $S$, on outer lines or external fields $F_i$. $G^{\mathrm{sw}}$ differs from $G$ only in the attachments at $\sigma_1,\sigma_2$, so $S$ has the same vertices, internal lines, differences, vertex localisations and anchor $\hat x$ in both, and the external lines of $S$ keep their indices, so that usefulness, the regrouped lines of the internal and external slices (which depend on the external indices only) and the transport of the multipliers to $\hat x$ (\cref{lem:a06d:S27}) agree. Hence the regrouped position kernel $K_S(y_1,y_2;\dots)$ of $S$, including its $R$-lines and the walled kernels, is the same function in both members, whatever the line kernels are. Its index structure is a matrix $M_S$ on the open index path from $\sigma_1$ to $\sigma_2$: every closed index cycle inside $S$ carries an even power of $q$ (else $\Tr q^{\mathrm{odd}}=0$ kills the term), and the total number of vector-leg ends at the vertices of $S$ is $2(\text{internal vector lines})+1$, odd, so the open path carries $q^{r}$ with $r$ odd times even polynomials in $q$, and $M_S^{\mathsf T}=-M_S$. At zeroth order the external fields are constants at $\hat x$ and the two local terms are $c(\hat x)F_1(\hat x)^{\mathsf T}M_SF_2(\hat x)$ and $c(\hat x)F_2(\hat x)^{\mathsf T}M_SF_1(\hat x)=-c(\hat x)F_1(\hat x)^{\mathsf T}M_SF_2(\hat x)$, with the same $c(\hat x)$; they cancel. Walled lines change nothing: they are scalar kernels times $\mathbf 1_N$ (\cref{lem:a06:WN1}(a)), so $M_S$ is unchanged; for slot fields which are kernels of several variables (the ends of $R^v_K$, $R^\gamma_K$), the zeroth-order term at $\hat x$ is their value at $\hat x$ in the slot variable, controlled by the norm \cite[(1.32)]{Ba3}. The remainders of the Taylor expansion are terms of the processing of Step~3(D$_4$), bounded term by term. This is Ba{\l}aban's reason for \cite[(2.18)]{Ba3}, ``the same expressions but with the opposite signs'' \cite[p.~444]{Ba3}, with the two members anchored at one vertex.

\emph{(b): the pinned root.} Here $u_1\ne u_2$ (the shapes T, B8, B10; in BL both scalar slots sit at $w$, so $u_1=u_2$ and (a) applies with the anchor at $w$). The local terms are $c_{S,1}(x)\Phi$ and $-c_{S,2}(x)\Phi$, with $c_{S,i}$ the regrouped coefficient with $u_i$ at $x$ and the anchor at $u_i$ (the sign from (a)). \emph{Split.} By the steps of \cref{thm:a06d:S48}: rescale to $\xi\Z^3$ ($\xi=L^{-j_0}$, $j_0$ the lowest external index), where by \cref{lem:a06:WN1}(f) and the rescaling of \cref{lem:a06:WN1}(a) every regrouped scalar line is $\breve G^v_l=C^\xi+R^v_l$ with a potential in $[0,\bar v]$ and $R^v_l$ as in \cref{lem:a06:WN1}(d), and every vector line is $C^\xi+R_l$ with the bulk remainder; move the multipliers to the anchor (\cref{lem:a06d:S27}); expand multilinearly, $c_{S,i}=c^C_{S,i}+c^R_{S,i}$. The shapes T, B8, B10 contain no \cite[(1.6)]{Ba3} or \cite[(1.7)]{Ba3} vertex, so no sharp indicator occurs and the internal sums are unrestricted on $\xi\Z^3$; and they contain no divergent proper subgraph (a two-scalar-leg or two-vector-leg subgraph inside a one-loop triangle or double line would need a bridge), so no inner counterterm occurs. \emph{The all-$C^\xi$ parts.} $c^C_{S,i}(x)$ is the same expression as in the bulk model: the same graph, vertex tensors and differences, the multipliers at the anchor, lines $C^\xi$, unrestricted sums. It is the zero-momentum value of one translation-invariant kernel on $\xi\Z^3$ with $u_i$ pinned: with $K^C_S(y_{u_1},y_{u_2},\dots)$ the all-$C^\xi$ kernel, $c^C_{S,1}(x)=\sum K^C_S(x,y_{u_2},\dots)=\sum K^C_S(0,y,\dots)$ and $c^C_{S,2}(x)=\sum K^C_S(y_{u_1},x,\dots)=\sum K^C_S(0,-y,\dots)$, both the total sum of $K^C_S(0,\cdot,\dots)$ over $\xi\Z^3$ (absolutely convergent at fixed $\xi$), so $c^C_{S,1}(x)=c^C_{S,2}(x)$, independent of $x$. The wall does not enter this term. \emph{The $R$-parts.} Each term of $c^R_{S,i}$ has at least one $R$-line. \Cref{lem:a06:WN1}(d) gives the bounds \eqref{eq:a06:WNR} for $R^v_l$ at every pair of points, with constant $C_R(1+\bar v)^3$, so $R^v_l$ is a remainder line of \cref{lem:a06d:S47}(b), and the step of \cref{thm:a06d:S48} for ``terms with at least one $R$-line, $D(S)=0$'' applies: the node is one-loop, its inner components are trees, the only dangerous ones are the \cite[(2.4)]{Ba3} subgraphs, treated by the summation by parts inside the coefficient (\cref{lem:a06d:S44}), every component after re-slicing has positive degree, and the bound function estimate \cite[(2.16)]{Ba3} bounds each $c^R_{S,i}$ by $C(1+\bar v)^p$ uniformly in $\xi$, $j_0$, $s$, $\eps$, the torus and $x$ (the kernel bounds are pointwise, and the anchor is fixed at $x$). The exponent $p$ counts the $R$-lines and the factors of $\bar v$ in their bounds, at most three lines here. \emph{Conclusion.} $|c_{S,1}(x)-c_{S,2}(x)|=|c^R_{S,1}(x)-c^R_{S,2}(x)|\le2C(1+\bar v)^p=: C^v_7$. This is the form of the $\phi\phi A$ clause of \cref{thm:a06d:S48}, $|c^{(0)}|\le C$, for the pair.

\emph{(c): bridged nodes.} Write $\Sigma_2(w,w')$ for the amputated kernel of $S_2$ between its leg point $w$ on the bridge and its other leg point $w'$, with $S_2$'s lines (walled or not, at the regrouped scale), and $F(y)\defeq(\partial^\xi_\mu\breve G_{l_0})(x_u,y)$ for the bridge factor, $u$'s difference acting on the bridge.

\emph{$u$ summed.} Choose the anchor $\hat x$ in $S_2$ (anchors of $\phi\phi A$ nodes are free, Step~3(D$_4$)). At zeroth order every external field of $S$ is a constant at $\hat x$, and so are $u$'s multipliers, which are smooth (\cref{lem:a06d:S27}; a one-vector-leg vertex carries no sharp indicator). The dependence of the regrouped kernel on $x_u$ is then the single factor $F$, and $\sum_{x_u}\xi^3(\partial^\xi_\mu\breve G_{l_0})(x_u,y)=0$ for every summable kernel: a forward difference in the summed variable telescopes, on $\xi\Z^3$ and on the torus. So $c=0$ for every member and any line kernels, before any $C^\xi/R$ split, and the same holds for the member with $S_2$ replaced by its counterterm vertex. No translation invariance of the lines is used.

\emph{$u$ pinned.} The anchor is $u$, at $x$. \emph{All-$C^\xi$ part.} On $\xi\Z^3$ with lines $C^\xi$ and unrestricted sums, $\sum_{w'}\Sigma^C_2(w,w')$ is a constant $g$ by translation invariance, and $\sum_w\xi^3(\partial^\xi_\mu C^\xi)(x,w)\,g=g\,\partial^\xi_\mu(C^\xi\mathbf 1)(x)=0$ since $C^\xi\mathbf 1\equiv1$. The counterterm member is the same computation with $\Sigma_2$ replaced by its own row sum. A sharp indicator inside $S_2$ occurs only at the \cite[(1.6)]{Ba3} tadpole, which cancels exactly against its counterterm member label by label ((Agr); the counterterm carries the same unit-cube label), so the pair has no indicator defect; the exchange bubble and the seagull loop carry only the smooth labels of their vector legs. So $c^C=0$ per member (and per pair). This is Ba{\l}aban's ``by translation invariance it can be written as a derivative of a constant'' \cite[p.~444]{Ba3} for \cite[(2.21c)]{Ba3}, and the exact cancellations he asserts for \cite[(2.21b), (2.21d), (2.21f), (2.21g)]{Ba3}, applied to the all-$C^\xi$ parts. \emph{Why the $R$-part is not bounded per member.} The term in which the bridge carries $R^v_{l_0}$ and $S_2$ carries lines $C^\xi$ equals $g\sum_w\xi^3(\partial^\xi_\mu R^v_{l_0})(x,w)=g\,\partial^\xi_\mu(R^v\mathbf 1)(x)$, where $g\sim\xi^{-1}$ (the node $S_2$ has $D(S_2)=-1$, and the $R$-line is not inside it, so it gains nothing for $S_2$). In the bulk $R\mathbf 1$ is constant ($-\Delta+m^2+aP$ maps constants to constants), and the term vanishes; under a wall $\breve G^v\mathbf 1=\mathbf 1-\breve G^v v$ is not constant, so the term is of order $\xi^{-1}$ and is not bounded uniformly in $\xi$. \emph{The pair.} Both members are in $\cC$ by (P). Let $f$ be the external field at $S_2$'s far slot $w'$ (constant at the anchor at zeroth order, Taylor-expanded at first order), and take the anchor convention (Agr) for $S_2$ with the anchor at the bridge end $w$, so that the counterterm member carries $\delta m^2_{S_2}(w)=\sum_{w'}\xi^3\Sigma_2(w,w')$, the row sum of the same kernel (the model's own counterterm, \cite[p.~417]{Ba3}). Then
\[
\sum_{w,w'}\xi^6F(w)\Sigma_2(w,w')f(w')-\sum_w\xi^3F(w)\,\delta m^2_{S_2}(w)f(w)\ =\ \sum_{w,w'}\xi^6F(w)\Sigma_2(w,w')\big[f(w')-f(w)\big],
\]
an identity for any kernels, with no symmetry used. Its zeroth-order term ($f$ constant) is $0$ exactly, including the terms in which an $R$-line sits inside $S_2$ (these have $D(S_2)+1=0$ and are paired with the counterterm member carrying the same $R$-line, as in the step of \cref{thm:a06d:S48} for $D(H)=-1$ components). The first-order term is $\sum_w\xi^3F(w)c_\mu(w)\,\partial_\mu f$ with $c_\mu(w)\defeq\sum_{w'}\xi^3\Sigma_2(w,w')(w'-w)_\mu$ the first moment of $S_2$; its all-$C^\xi$ part is bounded by \cite[(3.16)]{Ba3} (the cancellation $\sum_{y'}(\partial C^\xi\partial^*)(y'-y)=0$ of \cref{thm:a06d:S48}), and its $R^v$-part by $|\partial R^v\partial^*|\le C_R(1+\bar v)^3(\xi+\rho)^{-1}e^{-\delta_R\rho}$ (\cref{lem:a06:WN1}(d)) against the other line of the bubble, a sum convergent in $d=3$; so $|c_\mu(w)|\le C(1+\bar v)^3$, and $\sum_w\xi^3|F(w)|\le C(1+\bar v)^3$ by the bounds on $\partial C^\xi$ and $\partial R^v$. The H\"older remainder $f(w')-f(w)-(w'-w)\cdot\partial f(w)$ gives the node $S$ a differentiated external leg, after which every component has positive degree and \cite[(2.16)]{Ba3} bounds the sum. If the bridge attaches at $S_2$'s other leg and the anchor convention puts the counterterm at the far leg, the same identity holds with $f$ and $F$ exchanged, by the symmetry $\Sigma_2(w,w')=\Sigma_2(w',w)$ of the real symmetric scalar kernels at zero field (\cref{lem:a06:WN1}(a)). Altogether the $R$-part of the pair is $\le C(1+\bar v)^p=: C^v_{\mathrm{BR}}$, uniformly in $\xi$, $j_0$, $s$, $\eps$, the labels, the torus and $x$. The bridged shapes without a counterterm member (\cite[(2.21a), (2.21c), (2.21e)]{Ba3}: the one whose expression vanishes identically \cite[p.~442]{Ba3}, and those in which $S_2$ carries a difference on a leg or an extra power of $\eta$) need none: for them the all-$C^\xi$ part vanishes as above, and the $R$-part is convergent per member, since the difference on $S_2$'s leg, or the extra power of $\eta$, makes $D(S_2)\ge0$ before the gain of the $R$-line.
\end{proof}

\begin{remark}[reducible two-scalar-leg chains]\label{rem:a06:WNCH}
A one-particle-reducible two-scalar-leg chain of one-loop $D=-1$ parts is renormalised by the counterterms of its parts (Step~3(D$_4$)). The identity behind it, $T_{P_2}(b',c)=\delta_{b'c}\sum_{c'}K_{P_2}(b',c')$ with the same lines, indices and labels, holds for any kernels because the counterterm is computed with the model's own lines; so it holds verbatim in the walled model. The alternative route through translation invariance is not used.
\end{remark}

% a06_WN_d: part 4 of 4 of the walled reading of Ba3 Prop. 1 (input from sec/a06_BALtorus_d7b.tex): the proof of thm:a06:WN(ii), the constants, and two remarks on scope.

\begin{proof}[Proof of \cref{thm:a06:WN}\textup{(ii)}]
The proof of \cite[(1.33)]{Ba3} at $L=5$ for the bulk model, Steps~3(D$_3$)--3(D$_4$) with the twenty sketched steps listed after \cref{thm:a06d:S49}, is run for the walled model and for a class $\cC$ closed under the moves. It is organised as follows: first, where the scalar lines enter and what replaces each entry; second, where the class enters and what supplies each partner; third, the assembly and the induction.

\emph{Step 1: the scalar lines.} At zero background the scalar lines of every graph, slice, counterterm and cut kernel are zero-field kernels, and the proof uses them in exactly four ways. \emph{(1) Kernel bounds}: the group-(A) estimates (the slice bounds \cite[(2.6), (2.10)--(2.12)]{Ba3}, the kernel bounds of the cube step, the two-domain differences $\delta G_s(\square,\eta\Z^3,0)$, the cut-kernel estimate \cite[(3.1)]{Ba3} in the cluster form of \cref{cor:a06d:S24}, the Combes--Thomas bounds of the chains), each a composition of the pointwise bounds on the zero-field kernels with lattice sums; for walled lines these are \cref{lem:a06:WN1}(b),(c),(e), with constants polynomial in $\bar v$. The two-domain difference is bounded by the same resolvent identity as in the bulk, with \cref{lem:a06:WN1}(b) on each domain, since at zero field the bounds hold at every point of a Neumann domain and of the torus. \emph{(2) Algebra}: the slice identity, its rescaling, the gauge covariance used in the gauge-out of $B_0$, the cancellation of the operator in the differences of the $\tilde B'$-expansion and of the domain changes; for walled lines these are \cref{lem:a06:WN1}(a), and at zero background the $\tilde B'$-expansion, the constant field $B_0$ and its gauge-out are void. \emph{(3) The regrouped lines of the renormalisation}: in Step~3(D$_4$) the cumulative regrouped line of a node is, by the slice identity, the inverse of the operator pinned at the lowest external scale, which after rescaling to $\xi\Z^3$ is $\breve G_l=C^\xi+R_l$ (\cref{thm:a06d:S48}, \cref{lem:a06d:S47}); for walled lines the rescaled operator carries the potential $s^2v_{\mathrm{per}}(s\,\cdot)\in[0,\bar v]$ (\cref{lem:a06:WN1}(a),(f)), so $\breve G^v_l=C^\xi+R^v_l$ with the \emph{same} $C^\xi$ and with $R^v_l$ a remainder line of \cref{lem:a06d:S47}(b) with constant $C_R(1+\bar v)^3$ (\cref{lem:a06:WN1}(d)); hence every step of \cref{thm:a06d:S48} that concerns terms with at least one $R$-line (the one-loop $D=-1$ cancellation, which uses only the bounds on the $R$-line and the forward-difference structure of the $C^\xi$-line; the re-slicing of \cref{lem:a06d:S47}; the bound function estimate \cite[(2.16)]{Ba3}) runs verbatim, and every all-$C^\xi$ expression, with the inner counterterm members re-expanded with their own lines, is literally the same expression as in the bulk model \cite[pp.~437, 441]{Ba3}. \emph{(4) The counterterm values}: the model's own, by definition (the walled model). The torus-versus-$\xi\Z^3$ identification of \cite[pp.~418--419]{Ba3} is \cref{lem:a06:WN1}(f), with the image terms bounded, as in the bulk, by \cite[(1.33)]{Ba3} for the classes of lower order into which they decompose (the induction below). The sharp-indicator lemma \cref{lem:a06d:S46} uses only the decay of the internal lines, supplied by \cref{lem:a06:WN1}(b); the useless-counterterm lemma \cref{lem:a06d:S45} and the \cite[(2.4)]{Ba3} lemma \cref{lem:a06d:S44} are kernel-free given the bounds; \cref{lem:a06d:S41,lem:a06d:S42} are identities. Subgraphs with no scalar slot and an odd number of vector slots vanish at $B=0$ also with walled lines, which are scalar kernels times $\mathbf 1_N$ (\cref{lem:a06:WN1}(a)). A finite \cite[(1.7)]{Ba3} vertex with an $x$-dependent coefficient never sits inside a dangerous node (\cref{lem:a06d:S25}(b), \cite[Cor.~2.3]{Ba3}), so such coefficients enter only through their suprema, as the statement counts them.

\emph{Step 2: the class.} The class enters the bulk proof only at the dangerous nodes, as the sum over a node's partners with the rest of the graph fixed (\cref{lem:a06d:S41}(a): the counterterm pairing is an identity of class sums; \cref{thm:a06d:S48}: the local coefficient is the class sum at the node; \cref{thm:a06d:S28}), at the cut counterterms of partnered subgraphs (\cref{lem:a06d:S25}, which needs \cite[(1.33)]{Ba3} for the completed pieces of lower order), at the class-component cube (one label multiset), and in the regrouping after the cut (\cref{thm:a06d:S23}, kernel-free, (i)). Node by node, with walled lines: \emph{two-scalar-leg nodes} are paired with their counterterm member (P), the exact identities \cite[(3.19)]{Ba3} ($D=0$) and \cite[(3.11)]{Ba3} ($D=-1$) against the model's own counterterm with the same lines, the same localisation and (Agr); the $D=-1$ first moment has its all-$C^\xi$ part bounded by \cite[(3.16)]{Ba3} and its $R^v$-part by $|\partial R^v\partial^*|\le C(1+\bar v)^3(\xi+\rho)^{-1}e^{-\delta\rho}$, so $|c_\mu|\le C(1+\bar v)^3$; reducible chains by \cref{rem:a06:WNCH}. \emph{Two-vector-leg nodes}: the per-label two-vector-leg class at the node, which (W) supplies; the zeroth moment of its all-$C^\xi$ part vanishes by \cref{lem:a06d:S26}, applied as in \cref{thm:a06d:S48} to the full lines $C^\xi$ of the class sum (the lemma allows any $x$-dependent $|\phi|^2$ coefficients, so the inner counterterm and finite-vertex coefficients, and the potential itself, are admissible coefficients; \cref{lem:a06d:S26} is thus verbatim in the walled model), the first moment of its all-$C^\xi$ part by the reflection identity \cite[(3.30)]{Ba3} on $C^\xi$, which is model-free; the $R^v$-parts by the $R$-line steps of \cref{thm:a06d:S48}, with the zero-momentum remainder vertex kept. \emph{$\phi\phi A$ nodes}: in the bulk, \cref{thm:a06d:S48} bounds the zeroth moment by the reflection identity \cite[(3.37)]{Ba3} over the per-label full $\phi\phi A$ class; in a class closed under the moves this is replaced by (iii): at a bridgeless node the (X)-partner is in $\cC$ and (iii)(a) cancels the pair exactly at a common anchor, or, when the node carries the pinned root at one of its two slot vertices, (iii)(b) bounds the pair by $C^v_7$; at a bridged node no partner is needed, (iii)(c), the counterterm member (P) of $S_2$ being in $\cC$; and the pair bounds are of the form $|c^{(0)}|\le C$ that \cref{thm:a06d:S49} consumes for a $\phi\phi A$ local vertex. \emph{\cite[(2.4)]{Ba3} subgraphs}: summation by parts, \cref{lem:a06d:S44}, kernel-free. \emph{Vacuum components}: cancelled by the model's own $E_s[v]$ with the same labels and lines, up to image terms (Step~1). \emph{One- and three-vector-leg subgraphs}: vanish (Step~1). The cut counterterms of partnered subgraphs (\cref{lem:a06d:S25}) are classes of lower order of the completed pieces, which are closed under the moves (a move at a divergent subgraph of a piece is a move of the uncut member, and the completion of a partnered subgraph is by the moves); and the class-component cube depends on the label multiset only, common to $\cC$. Per-label full classes are closed under the moves, so the bulk statement for them is the special case $\cC=$ full class; and every partner that the walled node mechanisms use is in $\cC$ by the closure, while every other step is a bound per term.

\emph{Step 3: assembly and induction.} With Steps~1 and 2, \cref{thm:a06d:S48,thm:a06d:S49} hold for $\cC$ with walled lines, with the $\phi\phi A$ clause in the pair form of (iii); the induction on the number of lines inside \cref{thm:a06d:S48} and on the number of vertices in \cref{lem:a06d:S25} runs over classes closed under the moves of lower order (the completed pieces of a cut partnered subgraph, and the pieces into which an image term decomposes), to which the statement applies by the induction hypothesis; the base is the one-vertex graphs. The bound function estimate \cite[(2.16)]{Ba3} then gives \cite[(3.3)]{Ba3} for every class component, and the assembly of \cite[pp.~420--445]{Ba3} listed after \cref{thm:a06d:S49}, with the cluster cut and the cube steps, gives \cite[(1.33)]{Ba3} in the stated form, with $d_v$, $d_s$ counting the action couplings, the counterterm vertices as exact counterterms, and the finite vertices through their suprema (they have $D(v)=2$ whatever their coefficient, \cref{lem:a06d:S25}(b)); the localisation of an external field whose leg moves under the moves is (H$_A$); the slot fields of several variables are measured by \cite[(1.32)]{Ba3} in the slot variable, as the two-variable external fields of \cite[p.~419]{Ba3}; and in a two-point setting the root vertex is fixed at $x$ throughout, every kernel bound used being pointwise, so the bound is uniform in $x$.

\emph{Constants.} The walled inputs carry $(1+\bar v)^2$ (\cref{lem:a06:WN1}(b)), $(1+\bar v)^3$ (\cref{lem:a06:WN1}(d)) and the constants of \cref{lem:a06:WN1}(e), polynomial in $\bar v$; a graph of order $\le\bar n$ has a number of lines bounded by $\bar n$, so $O(1)^v$ is a polynomial in $\bar v$ of degree bounded in terms of $\bar n$, and $\delta^v_0$ is a fixed fraction of $\delta^w$, $\delta_R$ and the bulk rates. The parameter rows of Step~3(D) are kept, with $R_0$ replaced by $R^v_0\defeq R_0+10$ in every depth row for a scalar propagator bound and the torus side $\ge\max\{6\sqrt3\,r(t_s)+2,8/\delta^v_0\}$, all lower bounds on $R$, $r^\flat$ and the torus side or caps on $\eps_0$ of Ba{\l}aban's kind, fixed after $R_0$ and before $\eps_0$; nothing in the order of the choices changes. The constants are therefore independent of $\eps$, of $\eps_0\in(0,\eps^*_0]$ below the caps, of $s$, $K$, the labels, the torus side, $x$, and of $v$ beyond $\bar v$, and they depend on $L=5$, on the parameters of Step~3(D) and polynomially on $\bar v$.
\end{proof}

\begin{remark}[non-circularity of the uses in the rotation part]\label{rem:a06:WNcirc}
\Cref{thm:a06:WN}(ii) is a bound in which the suprema of the finite-vertex coefficients are parameters. Its uses in the rotation part at coupling order $4$ carry, as finite vertices, only the non-renormalising parts of the order-$2$ inner counterterms, whose sizes come from the order-$2$ statement, and the wall vertex, whose coefficient is $\le nt_s^2$; the sizes of the order-$4$ finite vertices are then \emph{consequences} of (ii) at order $4$, not inputs to it. The induction of the proof runs on the number of vertices inside the theorem.
\end{remark}

\begin{remark}[a bounded wall is not a boundary]\label{rem:a06:WNbdry}
\Cref{thm:a06:WN} is stated for a bounded site potential on the scalar lines with the vector lines unwalled, on a torus. It does not cover scalar or vector lines with Dirichlet or Neumann boundary conditions on a box, and the proof of (iii)(b) does not extend to them: a boundary changes the kernel by image terms $C^\xi(x-y^*)$ which, for $x$ and $y$ in the boundary layer, are as singular as $C^\xi$ itself and so are not remainder lines in the sense of \eqref{eq:a06:WNR}; with a Neumann vector line the odd part of the pinned pair is left uncancelled at every scale, and nothing bounds $|c_{S,1}-c_{S,2}|$ uniformly in $\xi$ in the layer. A bounded potential changes the kernel by $R^v$, which satisfies \eqref{eq:a06:WNR} at \emph{every} pair of points (\cref{lem:a06:WN1}(d); the multiplier $1-w$ is never differentiated, and no smoothness of $v$ is used), so there is no boundary layer, and the vector line is the translation-invariant torus propagator. This is why the rotation part compares King's cubes with walled tori and never applies the renormalisation theorem on a cube.
\end{remark}

